\documentclass[11pt,a4paper]{article}
\usepackage[utf8]{inputenc}
\usepackage[margin=1in]{geometry}
\usepackage{amsmath,amssymb}
\usepackage{booktabs}
\usepackage{xcolor}
\usepackage{listings}
\usepackage{hyperref}
\usepackage{titlesec}
\usepackage{microtype}
\usepackage{parskip}
\usepackage{array}
\usepackage{graphicx}

\definecolor{codebg}{RGB}{246,247,249}
\definecolor{kw}{RGB}{0,80,160}
\definecolor{cm}{RGB}{55,120,55}
\definecolor{st}{RGB}{160,40,40}
\definecolor{nu}{RGB}{130,55,170}
\definecolor{pp}{RGB}{100,50,0}
\definecolor{rulecol}{RGB}{185,192,202}

\lstdefinestyle{cpp20}{
  language=C++,
  backgroundcolor=\color{codebg},
  basicstyle=\ttfamily\footnotesize,
  keywordstyle=\color{kw}\bfseries,
  commentstyle=\color{cm}\itshape,
  stringstyle=\color{st},
  numberstyle=\color{nu}\tiny,
  emphstyle=\color{pp}\bfseries,
  emph={constexpr,noexcept,nodiscard,string_view,int32_t,size_t,
        ParticleID,DecayEvent,KineticEventRecord,EventKind,
        NuclearState,DaughterState,is_reserved,is_transportable,
        is_bookkeeping_only,is_decay_particle,is_energy_carrier,
        is_defect_or_virtual,is_ghost,is_atom_species,
        to_string,species_code_to_label},
  breaklines=true,
  frame=single,
  rulecolor=\color{rulecol},
  framesep=4pt,
  numbers=left,
  numbersep=8pt,
  tabsize=4,
  showstringspaces=false,
  xleftmargin=14pt,
  captionpos=b
}

\lstdefinestyle{clang}{
  language=C,
  backgroundcolor=\color{codebg},
  basicstyle=\ttfamily\footnotesize,
  keywordstyle=\color{kw}\bfseries,
  commentstyle=\color{cm}\itshape,
  stringstyle=\color{st},
  breaklines=true,
  frame=single,
  rulecolor=\color{rulecol},
  framesep=4pt,
  numbers=left,
  numbersep=8pt,
  tabsize=4,
  showstringspaces=false,
  xleftmargin=14pt,
  captionpos=b
}

\lstdefinestyle{tbl}{
  language={},
  backgroundcolor=\color{codebg},
  basicstyle=\ttfamily\scriptsize,
  breaklines=true,
  frame=single,
  rulecolor=\color{rulecol},
  framesep=4pt,
  numbers=none,
  showstringspaces=false,
  captionpos=b
}

\lstset{style=cpp20}

\titleformat{\section}{\large\bfseries}{}{0em}{}[\titlerule]
\titleformat{\subsection}{\normalsize\bfseries}{}{0em}{}

\begin{document}

\begin{center}
  {\LARGE\bfseries Unified Species Code Namespace}\\[4pt]
  {\large VSEPR-SIM -- Code Draft}\\[2pt]
  {\small\texttt{include/physics/particle\_id.hpp} \quad
          \texttt{legacy/c\_api/particle\_ids.h} \quad
          \texttt{include/physics/decay\_event\_builder.hpp}}\\[6pt]
  {\small v4.1.0 \quad June 2025}
\end{center}

\vspace{4pt}\hrule\vspace{8pt}
\tableofcontents
\newpage

%% ===================================================================
\section{Core Idea}
%% ===================================================================

Every row in the engine geometry table carries a single signed integer
called \texttt{species\_code}.  The sign tells you immediately what
kind of entity the coordinate belongs to:

\[
\texttt{species\_code} \in \mathbb{Z},\quad
\begin{cases}
\geq 1  & \text{standard atom template} \quad (Z = 1\ldots 118) \\
= 0     & \text{placeholder / unresolved} \\
\leq -1 & \text{reserved engine entity (negative class space)}
\end{cases}
\]

The negative range is a \emph{reserved class space}: every nonstandard,
derived, emitted, virtual, or bookkeeping entity that can appear in the
coordinate table gets a permanent negative code.  No separate type tag,
no string parsing, no secondary lookup.

\subsection{Geometry Table Integration}

A geometry table row has the shape:

\begin{lstlisting}[style=tbl,
  caption={Long-form geometry table with species\_code column}]
sample_id  bead_id  local_id  species_code  label        x        y        z       state
17         0        0          6             C            0.000    1.402    0.000   final
17         0        1          1             H            1.050    2.100    0.000   final
17         1        0         -1             alpha        2.500    0.300    0.200   emitted
17         1        1         -3             gamma        2.700    0.300    0.200   emitted
17         2        0         -4             eta          3.100    0.100    0.050   emitted
17         3        0         -17            ghost        0.000    0.000    0.000   virtual
\end{lstlisting}

\texttt{species\_code} is the master numeric identifier.
Positive or zero means standard atom/object template;
negative means reserved synthetic class.
The \texttt{label} column is derived from \texttt{species\_code}
via \texttt{species\_code\_to\_label()} and is never the source of truth.

%% ===================================================================
\section{The Reserved Negative Ladder}
%% ===================================================================

The negative range is partitioned into three consecutive bands:

\[
\underbrace{[-1,\ -8]}_{\substack{\text{physical emission /}\\\text{decay particles}}}
\quad
\underbrace{[-9,\ -14]}_{\substack{\text{energy / field /}\\\text{abstract carriers}}}
\quad
\underbrace{[-15,\ -18]}_{\substack{\text{defect / structural /}\\\text{bookkeeping}}}
\]

\bigskip
\begin{center}
\renewcommand{\arraystretch}{1.15}
\begin{tabular}{rllp{6.2cm}}
\toprule
\textbf{Code} & \textbf{Enum name} & \textbf{Label} & \textbf{Description} \\
\midrule
\multicolumn{4}{l}{\textit{Physical emission / decay particles}} \\
$-1$  & \texttt{alpha}            & \texttt{alpha}         & He-4 nucleus (alpha decay) \\
$-2$  & \texttt{beta\_minus}      & \texttt{beta-}         & Electron (beta-minus decay) \\
$-3$  & \texttt{gamma}            & \texttt{gamma}         & Gamma photon \\
$-4$  & \texttt{eta}              & \texttt{eta}           & QMD energy packet -- abstract
                                                             high-definition energy carrier;
                                                             extra to the standard kernel
                                                             energy unit \\
$-5$  & \texttt{antineutrino}     & \texttt{antineutrino}  & Antineutrino loss token \\
$-6$  & \texttt{neutrino}         & \texttt{neutrino}      & Neutrino loss token \\
$-7$  & \texttt{neutron\_free}    & \texttt{neutron}       & Free neutron \\
$-8$  & \texttt{proton\_free}     & \texttt{proton}        & Free proton \\
\midrule
\multicolumn{4}{l}{\textit{Energy / field / abstract carriers}} \\
$-9$  & \texttt{excitation\_packet}  & \texttt{excitation}    & Electronic excitation carrier \\
$-10$ & \texttt{ionization\_event}   & \texttt{ionization}    & Ionization bookkeeping token \\
$-11$ & \texttt{heat\_packet}        & \texttt{heat}          & Phonon / thermal energy packet \\
$-12$ & \texttt{charge\_cloud}       & \texttt{charge\_cloud} & Distributed charge carrier \\
$-13$ & \texttt{field\_source}       & \texttt{field\_source} & Field-origin sentinel \\
$-14$ & \texttt{field\_probe}        & \texttt{field\_probe}  & Field-sampling sentinel \\
\midrule
\multicolumn{4}{l}{\textit{Defect / structural / bookkeeping}} \\
$-15$ & \texttt{vacancy}            & \texttt{vacancy}           & Lattice vacancy defect \\
$-16$ & \texttt{interstitial}       & \texttt{interstitial}      & Lattice interstitial defect \\
$-17$ & \texttt{ghost}              & \texttt{ghost}             & Ghost particle -- virtual
                                                                    massless sentinel for field
                                                                    evaluation, constraint
                                                                    anchoring, or higher-order
                                                                    virtual entities \\
$-18$ & \texttt{transition\_state}  & \texttt{transition\_state} & Reaction coordinate marker \\
\bottomrule
\end{tabular}
\end{center}

\subsection{Notes on Selected Entries}

\paragraph{$-4$: eta packet.}
The \texttt{eta} entry is an abstract energy packet that lives above the
standard kernel energy unit.  It is specifically designed for
high-definition QMD simulations where sub-particle energy resolution is
required.  It can appear as an emitted row in the geometry table, carry
coordinates, and be tracked like any other entity.

\paragraph{$-17$: ghost particle.}
The ghost particle is massless and exerts no direct force.  Its coordinate
is meaningful: it marks a point in space used for field evaluation,
constraint anchoring, reaction coordinate sampling, or as a stand-in for
higher-order virtual entities.  A ghost row in the geometry table is valid
geometry; it never contributes to pair interactions.

%% ===================================================================
\section{\texttt{particle\_id.hpp} -- Full Listing}
%% ===================================================================

\begin{lstlisting}[style=cpp20,
  caption={\texttt{include/physics/particle\_id.hpp} -- enum class}]
namespace vsepr::physics {

enum class ParticleID : std::int32_t {
    placeholder            =  0,

    // Physical emission / decay particles
    alpha                  = -1,
    beta_minus             = -2,
    gamma                  = -3,
    eta                    = -4,   // QMD energy packet
    antineutrino           = -5,
    neutrino               = -6,
    neutron_free           = -7,
    proton_free            = -8,

    // Energy / field / abstract carriers
    excitation_packet      = -9,
    ionization_event       = -10,
    heat_packet            = -11,
    charge_cloud           = -12,
    field_source           = -13,
    field_probe            = -14,

    // Defect / structural / bookkeeping
    vacancy                = -15,
    interstitial           = -16,
    ghost                  = -17,
    transition_state       = -18
};

} // namespace vsepr::physics
\end{lstlisting}

\subsection{Sub-range Predicates}

\begin{lstlisting}[style=cpp20,
  caption={Constexpr sub-range predicates}]
// Full engine-reserved range [0, -99]
[[nodiscard]] constexpr bool is_reserved(ParticleID id) noexcept {
    const auto v = static_cast<std::int32_t>(id);
    return (v <= 0 && v >= -99);
}

// Physical decay / emission particles [-1, -8]
[[nodiscard]] constexpr bool is_decay_particle(ParticleID id) noexcept {
    const auto v = static_cast<std::int32_t>(id);
    return (v >= -8 && v <= -1);
}

// Physically propagates through matter
[[nodiscard]] constexpr bool is_transportable(ParticleID id) noexcept {
    switch (id) {
        case ParticleID::alpha:
        case ParticleID::beta_minus:
        case ParticleID::gamma:
        case ParticleID::neutron_free:
        case ParticleID::proton_free:
            return true;
        default: return false;
    }
}

// Never transported -- accounting token only
[[nodiscard]] constexpr bool is_bookkeeping_only(ParticleID id) noexcept {
    switch (id) {
        case ParticleID::eta:
        case ParticleID::antineutrino:
        case ParticleID::neutrino:
        case ParticleID::excitation_packet:
        case ParticleID::ionization_event:
        case ParticleID::heat_packet:
        case ParticleID::charge_cloud:
        case ParticleID::field_source:
        case ParticleID::field_probe:
        case ParticleID::transition_state:
            return true;
        default: return false;
    }
}

// Energy / field abstract carriers [-9, -14]
[[nodiscard]] constexpr bool is_energy_carrier(ParticleID id) noexcept {
    const auto v = static_cast<std::int32_t>(id);
    return (v >= -14 && v <= -9);
}

// Defect / structural / virtual entities [-15, -18]
[[nodiscard]] constexpr bool is_defect_or_virtual(ParticleID id) noexcept {
    const auto v = static_cast<std::int32_t>(id);
    return (v >= -18 && v <= -15);
}

// Ghost particle specifically
[[nodiscard]] constexpr bool is_ghost(ParticleID id) noexcept {
    return id == ParticleID::ghost;
}
\end{lstlisting}

\subsection{Label and \texttt{species\_code} Helpers}

\begin{lstlisting}[style=cpp20,
  caption={Geometry table label generation}]
[[nodiscard]] constexpr std::string_view to_string(ParticleID id) noexcept {
    switch (id) {
        case ParticleID::placeholder:        return "X";
        case ParticleID::alpha:              return "alpha";
        case ParticleID::beta_minus:         return "beta-";
        case ParticleID::gamma:              return "gamma";
        case ParticleID::eta:                return "eta";
        case ParticleID::antineutrino:       return "antineutrino";
        case ParticleID::neutrino:           return "neutrino";
        case ParticleID::neutron_free:       return "neutron";
        case ParticleID::proton_free:        return "proton";
        case ParticleID::excitation_packet:  return "excitation";
        case ParticleID::ionization_event:   return "ionization";
        case ParticleID::heat_packet:        return "heat";
        case ParticleID::charge_cloud:       return "charge_cloud";
        case ParticleID::field_source:       return "field_source";
        case ParticleID::field_probe:        return "field_probe";
        case ParticleID::vacancy:            return "vacancy";
        case ParticleID::interstitial:       return "interstitial";
        case ParticleID::ghost:              return "ghost";
        case ParticleID::transition_state:   return "transition_state";
        default:                             return "unknown";
    }
}

// Raw int32_t interface for geometry table I/O.
// Positive codes (atom Z) return "?" -- caller resolves via periodic table.
[[nodiscard]] constexpr std::string_view
species_code_to_label(std::int32_t code) noexcept {
    if (code > 0) return "?";
    return to_string(static_cast<ParticleID>(code));
}

// True if code is a standard atom template (Z in [1, 118])
[[nodiscard]] constexpr bool is_atom_species(std::int32_t code) noexcept {
    return (code >= 1 && code <= 118);
}

// True if code is a reserved engine entity (code <= 0)
[[nodiscard]] constexpr bool is_reserved_code(std::int32_t code) noexcept {
    return (code <= 0 && code >= -99);
}
\end{lstlisting}

%% ===================================================================
\section{Updated Decay Builders}
%% ===================================================================

The builder factories emit particles using the new enum names.
Beta-minus decay appends \texttt{antineutrino} (code $-5$) specifically
rather than a generic loss token.

\begin{lstlisting}[style=cpp20,
  caption={Alpha, beta-minus, and gamma decay factories}]
// S(Z,A) -> S(Z-2, A-4) + alpha
[[nodiscard]] inline DecayEvent make_alpha_decay(
    std::int32_t parent, std::int32_t daughter,
    double energy_eV, double time_s, double confidence
) {
    DecayEvent ev {};
    ev.parent_species_id   = parent;
    ev.daughter_species_id = daughter;
    ev.released_energy_eV  = energy_eV;
    ev.event_time_s        = time_s;
    ev.confidence          = confidence;
    ev.deposited_energy_fraction   = 0.85;
    ev.transported_energy_fraction = 0.15;
    ev.local_damage_score          = 0.90;
    (void)ev.add_emitted(ParticleID::alpha);         // -1
    return ev;
}

// S(Z,A) -> S(Z+1, A) + beta_minus + antineutrino
[[nodiscard]] inline DecayEvent make_beta_minus_decay(
    std::int32_t parent, std::int32_t daughter,
    double energy_eV, double time_s, double confidence
) {
    DecayEvent ev {};
    ev.parent_species_id   = parent;
    ev.daughter_species_id = daughter;
    ev.released_energy_eV  = energy_eV;
    ev.event_time_s        = time_s;
    ev.confidence          = confidence;
    ev.deposited_energy_fraction   = 0.35;
    ev.transported_energy_fraction = 0.65;
    ev.local_damage_score          = 0.30;
    (void)ev.add_emitted(ParticleID::beta_minus);    // -2
    (void)ev.add_emitted(ParticleID::antineutrino);  // -5
    return ev;
}

// S*(Z,A) -> S(Z,A) + gamma
[[nodiscard]] inline DecayEvent make_gamma_decay(
    std::int32_t parent, std::int32_t daughter,
    double energy_eV, double time_s, double confidence
) {
    DecayEvent ev {};
    ev.parent_species_id   = parent;
    ev.daughter_species_id = daughter;
    ev.released_energy_eV  = energy_eV;
    ev.event_time_s        = time_s;
    ev.confidence          = confidence;
    ev.deposited_energy_fraction   = 0.10;
    ev.transported_energy_fraction = 0.90;
    ev.local_damage_score          = 0.05;
    (void)ev.add_emitted(ParticleID::gamma);         // -3
    return ev;
}
\end{lstlisting}

%% ===================================================================
\section{Plain C Legacy API}
%% ===================================================================

\begin{lstlisting}[style=clang,
  caption={\texttt{legacy/c\_api/particle\_ids.h} -- updated C enum}]
typedef enum ParticleID {
    PARTICLE_PLACEHOLDER        =  0,

    /* Physical emission / decay particles */
    PARTICLE_ALPHA              = -1,
    PARTICLE_BETA_MINUS         = -2,
    PARTICLE_GAMMA              = -3,
    PARTICLE_ETA                = -4,   /* QMD energy packet */
    PARTICLE_ANTINEUTRINO       = -5,
    PARTICLE_NEUTRINO           = -6,
    PARTICLE_NEUTRON_FREE       = -7,
    PARTICLE_PROTON_FREE        = -8,

    /* Energy / field / abstract carriers */
    PARTICLE_EXCITATION         = -9,
    PARTICLE_IONIZATION         = -10,
    PARTICLE_HEAT_PACKET        = -11,
    PARTICLE_CHARGE_CLOUD       = -12,
    PARTICLE_FIELD_SOURCE       = -13,
    PARTICLE_FIELD_PROBE        = -14,

    /* Defect / structural / virtual */
    PARTICLE_VACANCY            = -15,
    PARTICLE_INTERSTITIAL       = -16,
    PARTICLE_GHOST              = -17,
    PARTICLE_TRANSITION_STATE   = -18
} ParticleID;

const char* particle_id_name(int id);
int particle_id_is_reserved(int id);
int particle_id_is_transportable(int id);
int particle_id_is_bookkeeping_only(int id);
int particle_id_is_atom_species(int id);
\end{lstlisting}

%% ===================================================================
\section{Query Classification Table}
%% ===================================================================

\begin{center}
\renewcommand{\arraystretch}{1.1}
\begin{tabular}{rlccccc}
\toprule
\textbf{Code} & \textbf{Name} &
  \rotatebox{70}{\texttt{reserved}} &
  \rotatebox{70}{\texttt{decay\_particle}} &
  \rotatebox{70}{\texttt{transportable}} &
  \rotatebox{70}{\texttt{energy\_carrier}} &
  \rotatebox{70}{\texttt{defect\_or\_virtual}} \\
\midrule
$-1$  & \texttt{alpha}              & $\bullet$ & $\bullet$ & $\bullet$ & & \\
$-2$  & \texttt{beta\_minus}        & $\bullet$ & $\bullet$ & $\bullet$ & & \\
$-3$  & \texttt{gamma}              & $\bullet$ & $\bullet$ & $\bullet$ & & \\
$-4$  & \texttt{eta}                & $\bullet$ & $\bullet$ & & & \\
$-5$  & \texttt{antineutrino}       & $\bullet$ & $\bullet$ & & & \\
$-6$  & \texttt{neutrino}           & $\bullet$ & $\bullet$ & & & \\
$-7$  & \texttt{neutron\_free}      & $\bullet$ & $\bullet$ & $\bullet$ & & \\
$-8$  & \texttt{proton\_free}       & $\bullet$ & $\bullet$ & $\bullet$ & & \\
$-9$  & \texttt{excitation\_packet} & $\bullet$ & & & $\bullet$ & \\
$-10$ & \texttt{ionization\_event}  & $\bullet$ & & & $\bullet$ & \\
$-11$ & \texttt{heat\_packet}       & $\bullet$ & & & $\bullet$ & \\
$-12$ & \texttt{charge\_cloud}      & $\bullet$ & & & $\bullet$ & \\
$-13$ & \texttt{field\_source}      & $\bullet$ & & & $\bullet$ & \\
$-14$ & \texttt{field\_probe}       & $\bullet$ & & & $\bullet$ & \\
$-15$ & \texttt{vacancy}            & $\bullet$ & & & & $\bullet$ \\
$-16$ & \texttt{interstitial}       & $\bullet$ & & & & $\bullet$ \\
$-17$ & \texttt{ghost}              & $\bullet$ & & & & $\bullet$ \\
$-18$ & \texttt{transition\_state}  & $\bullet$ & & & & $\bullet$ \\
\bottomrule
\end{tabular}
\end{center}

Note: \texttt{is\_bookkeeping\_only} returns true for all entries where
\texttt{is\_transportable} is false and \texttt{is\_reserved} is true,
i.e.\ codes $-4$ through $-6$, codes $-9$ through $-14$, and $-18$.

%% ===================================================================
\section{Geometry Table Dispatch Pattern}
%% ===================================================================

\subsection{Writing a Row}

\begin{lstlisting}[style=cpp20,
  caption={Writing a geometry table row for any species code}]
void write_row(std::ostream& out,
               int sample_id, int bead_id, int local_id,
               std::int32_t species_code,
               double x, double y, double z,
               std::string_view state)
{
    using namespace vsepr::physics;

    std::string_view label;
    if (is_atom_species(species_code)) {
        label = element_symbol(species_code); // periodic table
    } else {
        label = species_code_to_label(species_code);
    }

    out << sample_id    << '\t'
        << bead_id      << '\t'
        << local_id     << '\t'
        << species_code << '\t'
        << label        << '\t'
        << x << '\t' << y << '\t' << z << '\t'
        << state        << '\n';
}

// Examples:
write_row(out, 17, 0, 0,  6,  x, y, z, "final");   // Carbon
write_row(out, 17, 1, 0, -1,  x, y, z, "emitted");  // alpha
write_row(out, 17, 2, 0, -4,  x, y, z, "emitted");  // eta
write_row(out, 17, 3, 0, -17, x, y, z, "virtual");  // ghost
\end{lstlisting}

\subsection{Reading and Dispatching}

\begin{lstlisting}[style=cpp20,
  caption={Dispatch on species\_code while reading geometry table}]
void dispatch_row(std::int32_t sc, double x, double y, double z)
{
    using namespace vsepr::physics;
    const auto id = static_cast<ParticleID>(sc);

    if (is_atom_species(sc)) {
        add_atom(sc, x, y, z);               // force-field
    }
    else if (is_decay_particle(id)) {
        register_emitted(sc, x, y, z);        // transport
    }
    else if (is_energy_carrier(id)) {
        register_energy_token(sc, x, y, z);   // accounting
    }
    else if (is_ghost(id)) {
        register_field_point(x, y, z);        // no force
    }
    else if (is_defect_or_virtual(id)) {
        register_defect(sc, x, y, z);         // structural
    }
}
\end{lstlisting}

%% ===================================================================
\section{Relation to Existing Infrastructure}
%% ===================================================================

\begin{center}
\begin{tabular}{lll}
\toprule
\textbf{Component} & \textbf{Location} & \textbf{Role} \\
\midrule
\texttt{vsepr::nuclear::DecayMode}       & \texttt{decay\_chains.hpp}      & How a chain step progresses \\
\texttt{atomistic::nuclear::DecayType}   & \texttt{nuclear\_stability.hpp} & Dominant mode prediction for $Z$ \\
\texttt{vsepr::physics::ParticleID}      & \texttt{particle\_id.hpp}       & Integer code in the geometry row \\
\texttt{XYZAtom::atom\_type}             & \texttt{xyz\_format.hpp}        & Force-field type (always positive) \\
\bottomrule
\end{tabular}
\end{center}

\texttt{DecayMode} and \texttt{DecayType} answer ``what kind of
transition occurred?''  \texttt{ParticleID} / \texttt{species\_code}
answers ``what integer code is in this coordinate row?''
The four components are orthogonal.

\texttt{XYZAtom::atom\_type} carries force-field type codes and is
always positive.  \texttt{species\_code} is a separate higher-level
column; negative values there do not conflict with \texttt{atom\_type}.

\end{document}
